\section{Sign Conventions and Units}
\label{sec:sign_conventions_and_units}

Relativistic formulas are sensitive to conventions. Two references may describe the same physics while displaying different signs or factors of \(c\). A calculation is reliable only when its metric signature, coordinate convention, index placement, and unit system are stated and used consistently.

This section collects the conventions adopted in these notes, compares them with common alternatives, and provides a systematic procedure for translating formulas.

\subsection{Conventions used in these notes}

The spacetime coordinates are

\[
\boxed{
x^{\mu}=(ct,x,y,z)
}.
\]

Greek indices run over

\[
\mu,\nu,\alpha,\beta=0,1,2,3,
\]

while Latin spatial indices run over

\[
i,j,k=1,2,3.
\]

The Minkowski metric is

\[
\boxed{
\eta_{\mu\nu}
=
\eta^{\mu\nu}
=
\operatorname{diag}(1,-1,-1,-1)
}.
\]

This is the mostly-minus or time-positive signature.

The interval is

\[
\boxed{
\mathrm{d}s^{2}
=
c^{2}\mathrm{d}t^{2}
-
\mathrm{d}\mathbf{x}^{2}
}.
\]

Timelike vectors have positive squared norm, null vectors have zero squared norm, and spacelike vectors have negative squared norm.

\subsection{Raising and lowering signs}

For a contravariant four-vector

\[
A^{\mu}=(A^{0},\mathbf{A}),
\]

the corresponding covector is

\[
\boxed{
A_{\mu}=(A^{0},-\mathbf{A})
}.
\]

Thus

\[
A_{0}=A^{0},
\qquad
A_{i}=-A^{i}.
\]

Every raised or lowered spatial index contributes a minus sign in Cartesian Minkowski coordinates.

For a rank-two tensor, the number of sign changes depends on the number of spatial indices. For example,

\[
T_{00}=T^{00},
\]

\[
T_{0i}=-T^{0i},
\]

and

\[
T_{ij}=T^{ij}.
\]

The last equality contains two spatial sign changes, which cancel.

\subsection{Derivative conventions}

Because

\[
x^{0}=ct,
\]

the covariant four-gradient is

\[
\boxed{
\partial_{\mu}
=
\left(
\frac{1}{c}\partial_{t},
\nabla
\right)
}.
\]

Raising the index gives

\[
\boxed{
\partial^{\mu}
=
\left(
\frac{1}{c}\partial_{t},
-\nabla
\right)
}.
\]

The d'Alembertian is

\[
\boxed{
\Box
=
\partial_{\mu}\partial^{\mu}
=
\frac{1}{c^{2}}\partial_{t}^{2}
-
\nabla^{2}
}.
\]

These formulas are fixed by the coordinate choice and metric signature. They should not be memorized independently of those conventions.

\subsection{The opposite metric signature}

Many texts use

\[
\widetilde{\eta}_{\mu\nu}
=
\operatorname{diag}(-1,+1,+1,+1).
\]

Then the interval is

\[
\mathrm{d}\widetilde{s}^{2}
=
-c^{2}\mathrm{d}t^{2}
+
\mathrm{d}\mathbf{x}^{2}.
\]

Timelike vectors have negative norm and spacelike vectors have positive norm. The raised derivative becomes

\[
\widetilde{\partial}^{\mu}
=
\left(
-\frac{1}{c}\partial_{t},
\nabla
\right),
\]

and the contracted operator is

\[
\widetilde{\partial}_{\mu}
\widetilde{\partial}^{\mu}
=
-\frac{1}{c^{2}}\partial_{t}^{2}
+
\nabla^{2}.
\]

Thus

\[
\widetilde{\Box}=-\Box
\]

if both authors define the d'Alembertian directly as the metric contraction.

The two signatures describe the same physics. Problems arise only when formulas from different conventions are combined without translation.

\subsection{Translating a scalar-field kinetic term}

With the mostly-minus signature, the Lorentz scalar is

\[
\partial_{\mu}\phi\partial^{\mu}\phi
=
\frac{1}{c^{2}}\phi_{t}^{2}
-
|\nabla\phi|^{2}.
\]

A standard scalar-field Lagrangian density therefore contains

\[
\mathcal{L}_{\mathrm{kin}}
=
\frac{1}{2}
\partial_{\mu}\phi\partial^{\mu}\phi.
\]

With the mostly-plus signature, the contraction has the opposite sign:

\[
\widetilde{\partial}_{\mu}\phi
\widetilde{\partial}^{\mu}\phi
=
-rac{1}{c^{2}}\phi_{t}^{2}
+
|\nabla\phi|^{2}.
\]

To preserve a positive time-derivative contribution in the Lagrangian, one commonly writes

\[
\widetilde{\mathcal{L}}_{\mathrm{kin}}
=
-\frac{1}{2}
\widetilde{\partial}_{\mu}\phi
\widetilde{\partial}^{\mu}\phi.
\]

The difference is conventional. The expanded physical kinetic and gradient terms agree.

\subsection{The position of indices is part of the formula}

The expressions

\[
A^{\mu},
\qquad
A_{\mu}
\]

represent related but different component arrays. Likewise,

\[
\partial_{\mu}
\qquad\text{and}\qquad
\partial^{\mu}
\]

are not interchangeable.

An equation such as

\[
A^{\mu}=B_{\mu}
\]

is not a valid tensor equation merely because both sides have four components. Their transformation types differ.

A valid equality must have matching free indices in matching positions, for example

\[
A^{\mu}=B^{\mu}
\]

or

\[
A_{\mu}=B_{\mu}.
\]

If the metric is used explicitly, one may write

\[
A_{\mu}=\eta_{\mu\nu}A^{\nu}.
\]

\subsection{Explicit factors of \(c\)}

Keeping \(c\) explicit clarifies dimensions. Important formulas include

\[
x^{0}=ct,
\]

\[
U^{\mu}=\gamma(c,\mathbf{v}),
\]

\[
p^{\mu}
=
\left(
\frac{E}{c},
\mathbf{p}
\right),
\]

\[
E^{2}
=
|\mathbf{p}|^{2}c^{2}
+
m^{2}c^{4},
\]

and

\[
\Box
=
\frac{1}{c^{2}}\partial_{t}^{2}
-
\nabla^{2}.
\]

The factor \(c\) converts time units to length units and energy units to momentum-times-speed units.

\subsection{Natural units with \(c=1\)}

In relativistic theory, it is common to choose units in which

\[
\boxed{c=1}.
\]

Then time and length have the same units, velocity is dimensionless, and the coordinates become

\[
x^{\mu}=(t,\mathbf{x}).
\]

The four-velocity becomes

\[
U^{\mu}=\gamma(1,\mathbf{v}),
\]

and the four-momentum becomes

\[
p^{\mu}=(E,\mathbf{p}).
\]

The energy-momentum relation simplifies to

\[
\boxed{
E^{2}=|\mathbf{p}|^{2}+m^{2}
}.
\]

The d'Alembertian becomes

\[
\boxed{
\Box=\partial_{t}^{2}-\nabla^{2}
}.
\]

These formulas are compact, but the dimensions have not disappeared. The unit convention identifies quantities whose SI dimensions differ by powers of \(c\).

\subsection{Restoring powers of \(c\)}

A formula written with \(c=1\) can be translated back by dimensional analysis.

For example,

\[
E^{2}=p^{2}+m^{2}
\]

cannot be an SI equation because its three terms have different SI dimensions. The required factors are

\[
E^{2}
=
p^{2}c^{2}
+
m^{2}c^{4}.
\]

Likewise,

\[
x^{0}=t
\]

in natural units becomes

\[
x^{0}=ct
\]

when all coordinates are assigned dimensions of length.

A reliable restoration procedure is:

\begin{enumerate}
    \item identify the target physical dimension of every term;
    \item insert powers of \(c\) until all terms match;
    \item check a familiar limit or special case;
    \item compare with the definition of the corresponding four-vector.
\end{enumerate}

\subsection{Units with \(\hbar=1\)}

Quantum and high-energy field theory often also uses

\[
\hbar=1.
\]

Then energy, inverse time, inverse length, mass, and momentum are related by the same unit scale. For example,

\[
E=\hbar\omega
\]

becomes

\[
E=\omega,
\]

and

\[
\mathbf{p}=\hbar\mathbf{k}
\]

becomes

\[
\mathbf{p}=\mathbf{k}.
\]

The present classical course will usually keep \(c\) explicit in derivations. Later chapters may occasionally mention natural-unit forms, but the convention will be stated when used.

\subsection{Dimensions of common relativistic objects}

With \(x^{0}=ct\),

\[
[x^{\mu}]=\mathrm{m}.
\]

Therefore

\[
[\partial_{\mu}]=\mathrm{m}^{-1},
\]

and

\[
[\Box]=\mathrm{m}^{-2}.
\]

The four-velocity has dimensions

\[
[U^{\mu}]=\mathrm{m\,s^{-1}},
\]

because proper time is measured in seconds.

The four-momentum components all have momentum dimensions:

\[
[p^{\mu}]=\mathrm{kg\,m\,s^{-1}}.
\]

Indeed,

\[
\left[\frac{E}{c}\right]
=
\frac{\mathrm{kg\,m^{2}\,s^{-2}}}
{\mathrm{m\,s^{-1}}}
=
\mathrm{kg\,m\,s^{-1}}.
\]

\subsection{Spatial-index conventions}

When spatial indices are used, the Euclidean Kronecker delta

\[
\delta_{ij}
\]

is often used to contract ordinary three-vectors:

\[
|\mathbf{A}|^{2}
=
\delta_{ij}A^{i}A^{j}.
\]

The spatial part of the Minkowski metric is

\[
\eta_{ij}=-\delta_{ij}.
\]

Therefore

\[
\eta_{ij}A^{i}A^{j}
=
-|\mathbf{A}|^{2}.
\]

One must distinguish a three-dimensional Euclidean contraction from a four-dimensional Minkowski contraction.

\subsection{Levi--Civita conventions}

Later electromagnetic calculations may use the antisymmetric Levi--Civita symbol. A convention must specify the sign of

\[
\epsilon^{0123}.
\]

With the mostly-minus metric, lowering all four indices changes the sign because

\[
\det(\eta)=-1.
\]

Thus, if

\[
\epsilon^{0123}=+1,
\]

then

\[
\epsilon_{0123}=-1.
\]

Whenever the Levi--Civita symbol appears later, this sign convention will be stated explicitly rather than assumed silently.

\subsection{A convention-translation example}

Suppose one reference using the mostly-minus signature writes

\[
\Box\phi+m^{2}\phi=0
\]

in units with \(c=1\).

Another reference uses the mostly-plus signature and defines

\[
\widetilde{\Box}
=
\widetilde{\partial}_{\mu}
\widetilde{\partial}^{\mu}
=-\Box.
\]

The same equation is then written as

\[
-\widetilde{\Box}\phi+m^{2}\phi=0,
\]

or equivalently

\[
\widetilde{\Box}\phi-m^{2}\phi=0.
\]

The apparent sign disagreement disappears once the definitions are translated consistently.

\subsection{A diagnostic checklist}

Before accepting a relativistic formula, ask:

\begin{enumerate}
    \item Which metric signature is being used?
    \item Is the zeroth coordinate \(t\) or \(ct\)?
    \item Are factors of \(c\) explicit or has \(c=1\) been adopted?
    \item Are upper and lower indices distinguished?
    \item Does every contraction pair one upper and one lower index?
    \item What is the author's definition of \(\Box\)?
    \item Are spatial contractions Euclidean or inherited from the Minkowski metric?
    \item Do all terms have consistent physical dimensions?
\end{enumerate}

Most sign disputes can be resolved by answering these questions before manipulating the formula.

\subsection{Common mistakes}

Typical errors include:

\begin{enumerate}
    \item copying a formula from the opposite signature without translating it;
    \item forgetting that lowering a spatial index changes its sign;
    \item mixing \(x^{0}=t\) with formulas derived for \(x^{0}=ct\);
    \item setting \(c=1\) in only part of an equation;
    \item restoring powers of \(c\) by memory rather than dimensional analysis;
    \item treating \(\delta_{ij}\) and \(\eta_{ij}\) as the same object;
    \item assuming all authors use the same d'Alembertian convention;
    \item declaring two formulas inconsistent before comparing their conventions.
\end{enumerate}

\subsection{Chapter summary}

This chapter established the geometric and notational language required for relativistic field theory.

An event is represented in an inertial frame by

\[
x^{\mu}=(ct,\mathbf{x}),
\]

and the Minkowski metric

\[
\eta_{\mu\nu}
=
\operatorname{diag}(1,-1,-1,-1)
\]

defines the invariant interval

\[
\mathrm{d}s^{2}
=
c^{2}\mathrm{d}t^{2}
-
\mathrm{d}\mathbf{x}^{2}.
\]

Four-vectors transform as

\[
A'^{\mu}
=
\Lambda^{\mu}{}_{\nu}A^{\nu},
\]

while the metric raises and lowers indices. Complete contractions such as

\[
A_{\mu}B^{\mu}
\]

are Lorentz scalars.

The standard boost mixes time and the longitudinal spatial coordinate while preserving the interval. Four-dimensional differentiation is organized through

\[
\partial_{\mu}
=
\left(
\frac{1}{c}\partial_{t},
\nabla
\right),
\]

and

\[
\Box
=
\partial_{\mu}\partial^{\mu}
=
\frac{1}{c^{2}}\partial_{t}^{2}
-
\nabla^{2}.
\]

The principal invariants include

\[
\Delta x^{2},
\qquad
U^{2}=c^{2},
\qquad
p^{2}=m^{2}c^{2},
\qquad
k_{\mu}x^{\mu}.
\]

The chapter also established a working discipline: every relativistic calculation must state its signature, coordinate convention, index rules, and units before signs are compared.

The next chapter applies this formalism to the real scalar field. Lorentz-invariant derivative and potential terms will be assembled into an action, and the massless and massive Klein--Gordon equations will be derived and analyzed.